\chapter{1952年全国卷}

\section{填空题}


\begin{problem}
分解因式：\(x^{4} - y^{4} =\)\fillinblank{}.
\end{problem}

\begin{answer}
\(\left( x-y \right)\left( x+y \right)\left( x^{2}+y^{2} \right)\).
\end{answer}

\begin{problem}
若 \(\log_{10} 2x = 2\log_{10} x\)，问 \(x=\)\fillinblank{}.
\end{problem}

\begin{answer}
\(2\).
\end{answer}

\begin{problem}
若方程式 \(x^{3} + bx^{2} + cx + d = 0\) 之三根为 \(1\)，\(-1\)，\(\frac{1}{2}\)，则 \(c=\)\fillinblank{}.
\end{problem}

\begin{answer}
\(-1\).
\end{answer}

\begin{problem}
若 \(\sqrt{x^{2}+7}-4=0\)，则 \(x=\)\fillinblank{}.
\end{problem}

\begin{answer}
\(\pm 3\).
\end{answer}

\begin{problem}
\(\begin{vmatrix}1 & 2 & 3 \\ 4 & 5 & 0 \\ 3 & 2 & 1\end{vmatrix}=\)\fillinblank{}.
\end{problem}

\begin{answer}
\(-24\).
\end{answer}

\begin{problem}
两个圆的半径都是 \(4\text{ 寸}\)，并且一个圆通过另一圆的圆心，则这两个圆的公共弦之长是\fillinblank{}\(\text{寸}\).
\end{problem}

\begin{answer}
\(4\sqrt{3}\).
\end{answer}

\begin{problem}
三角形 \(\triangle ABC\) 的面积是 \(60\text{ 平方寸}\)，\(M\) 是 \(AB\) 的中点，\(N\) 是 \(AC\) 的中点，则 \(\triangle AMN\) 的面积是\fillinblank{}\(\text{平方寸}\).
\end{problem}

\begin{answer}
\(15\).
\end{answer}

\begin{problem}
正十边形的一内角是\fillinblank{}\(\text{度}\).
\end{problem}

\begin{answer}
\(144^{\circ}\).
\end{answer}

\begin{problem}
祖冲之的圆周率 \(\pi=\)\fillinblank{}.
\end{problem}

\begin{answer}
\(\frac{22}{7}\)，\(\frac{355}{113}\)，\(3.14159265\).
\end{answer}

\begin{problem}
球的面积等于大圆面积的\fillinblank{} 倍.
\end{problem}

\begin{answer}
\(4\).
\end{answer}

\begin{problem}
直圆锥之底之半径为 \(3\text{ 尺}\)，斜高为 \(5\text{ 尺}\)，则其体积为\fillinblank{}\(\text{立方尺}\).
\end{problem}

\begin{answer}
\(12\pi\).
\end{answer}

\begin{problem}
正多面体有\fillinblank{} 种，其名称为\fillinblank{}.
\end{problem}

\begin{answer}
\(5\)，正四面体，正六面体，正八面体，正十二面体，正二十面体.
\end{answer}

\begin{problem}
已知 \(\sin \theta = \frac{1}{3}\)，求 \(\cos 2\theta=\)\fillinblank{}.
\end{problem}

\begin{answer}
\(\frac{7}{9}\).
\end{answer}

\begin{problem}
方程式 \(\tan 2x = 1\) 的通解为 \(x=\)\fillinblank{}.
\end{problem}

\begin{answer}
\(\frac{1}{2}\left( n\pi + \frac{\pi}{4} \right)\).
\end{answer}

\begin{problem}
太阳仰角为 \(30^{\circ}\) 时塔影长 \(5\text{ 丈}\)，求塔高 \(=\)\fillinblank{}.
\end{problem}

\begin{answer}
\(\frac{5}{3}\sqrt{3}\).
\end{answer}

\begin{problem}
三角形 \(\triangle ABC\) 之 \(b\) 边为 \(3\text{ 寸}\)，\(c\) 边为 \(4\text{ 寸}\)，\(A\) 角为 \(30^{\circ}\)，则 \(\triangle ABC\) 的面积为\fillinblank{}\(\text{平方寸}\).
\end{problem}

\begin{answer}
\(3\).
\end{answer}

\begin{problem}
已知一直线经过点 \(\left( 2,-3 \right)\)，其斜率为 \(-1\)，则此直线之方程式为\fillinblank{}.
\end{problem}

\begin{answer}
\(x + y + 1 = 0\).
\end{answer}

\begin{problem}
若原点在一圆上，而此圆的圆心为点 \(\left( 3,4 \right)\)，则此圆的方程式为\fillinblank{}.
\end{problem}

\begin{answer}
\(x^{2} + y^{2} - 6x - 8y = 0\).
\end{answer}

\begin{problem}
原点至 \(3x + 4y + 1 = 0\) 之距离 \(=\)\fillinblank{}.
\end{problem}

\begin{answer}
\(\frac{1}{5}\).
\end{answer}

\begin{problem}
抛物线 \(y^{2} - 8x + 6y + 17 = 0\) 之顶点之坐标为\fillinblank{}.
\end{problem}

\begin{answer}
\(\left( 1,-3 \right)\).
\end{answer}

\section{解答题}


\begin{problem}
解方程式 \(x^{4} + 5x^{3} - 7x^{2} - 8x - 12 = 0\).
\end{problem}

\begin{solution}
\(2\)，\(-6\)，\(\omega\)，\(\omega^{2}\).
\end{solution}

\begin{problem}
\(\triangle ABC\) 中，\(\angle A\) 的外分角线与此三角形的外接圆相交于 \(D\)，求证：\(BD = CD\).
\end{problem}

\begin{problem}
设三角形的边长为 \(a = 4\)，\(b = 5\)，\(c = 6\)，其对角依次为 \(A\)，\(B\)，\(C\).

\begin{enumerate}
    \item 求 \(\cos C\).
    \item 求 \(\sin C\)，\(\sin B\)，\(\sin A\).
    \item 问 \(A\)，\(B\)，\(C\) 三个角各为锐角或钝角？
\end{enumerate}
\end{problem}

\begin{solution}
\begin{enumerate}
    \item \(\cos C = \frac{1}{8}\).
    \item \(\sin C = \frac{3}{8}\sqrt{7}\)，\(\sin B = \frac{5}{16}\sqrt{7}\)，\(\sin A = \frac{1}{4}\sqrt{7}\).
    \item \(A\)，\(B\)，\(C\) 皆为锐角.
\end{enumerate}
\end{solution}

\begin{problem}
一椭圆通过 \(\left( 2,3 \right)\) 及 \(\left( -1,4 \right)\) 两点，中心为原点，长短轴重合于坐标轴，试求其长短轴及焦点.
\end{problem}

\begin{solution}
长轴：\(\frac{2}{3}\sqrt{165}\)，短轴：\(\frac{2}{7}\sqrt{385}\)，焦点：\(\left( 0,\pm \frac{2}{21}\sqrt{1155} \right)\).
\end{solution}
